% ---------------------------------------------------------------------------------
% 1. INTRODUCTION
% Guide: context and objective; dataset (5,861 images 320x320, 3,516/1,172/1,173,
% class distribution, original vs input size); scope of the work.
% ---------------------------------------------------------------------------------
\section{Introduction}\label{sec:intro}

Given a single RGB satellite image, the task is to predict one of seven FireRisk classes: Water, Very Low, Low, Moderate, High, Very High and Non-burnable. It is an image-classification problem, since the label describes the fire risk of the area rather than the presence of fire. Two properties make it difficult. First, the classes are strongly imbalanced: Water accounts for only 2.4\% of the images, whereas Very Low accounts for 31\%. Second, the risk levels are ordinal and visually similar, so neighbouring levels such as Low, Moderate and High are easily confused.

We work with a subset of FireRisk \citep{shen2023firerisk} containing 5,861 images of $320\times320$ pixels, split in a stratified way into 3,516 training, 1,172 validation and 1,173 test images whose labels are hidden (\cref{tab:data}). Both final networks receive $224\times224$ crops of these images.

We developed two systems. The CUSTOM system is a VGG-style CNN trained from scratch and tuned with a staged random search, while the FINE-TUNING (FT) system adapts ImageNet-pretrained torchvision models, among which Swin-T was selected. Along the way we studied the effect of network width and depth, training length, class weighting, augmentation and image mixing, ensembles, and retraining on train+val.
